\documentclass{article}
\usepackage[utf8]{inputenc}
\usepackage{url}
\usepackage{hyperref}
\usepackage{anysize} %Margen
\usepackage{graphicx} %imagenes
\usepackage[spanish]{babel}
\usepackage{sectsty}
\usepackage{hyperref}
\usepackage{float}

\begin{document}
\marginsize{2cm}{2cm}{1cm}{2cm} 

\begin{center}
  {\LARGE \scshape Primer Proyecto de Software\\\vspace{10mm} }
  \rule{0.8\textwidth}{.8pt}\\
\end{center}

\section*{Involucrados}
\subsection*{Nombre del equipo} \textit{\textbf{ILAKECH ALAKECH}}
\subsection*{Logo del equipo}
\begin{center}
  \includegraphics[scale=.2]{../imagenes/logo.jpg}
\end{center}

\subsection*{Equipo de trabajo}
\begin{itemize}
\item Galeana Araujo Emiliano. Líder de equipo.
\item Macías Gómez Jorge. Técnico.
\item Torres Bucio Miriam. Responsable de Calidad.
\item Segoviano Gómez Víctor Manuel. Responsable de Colaboración.
\item Velez Cardenas Juan Luis Guillermo. Responsable de Colaboración.
\end{itemize}

\subsection*{Docente} Hanna Jadwiga Oktaba\\
\rule{0.8\textwidth}{.8pt}\\

\section*{Nombre del producto de software}
PUMA Eventos\\
\rule{0.8\textwidth}{.8pt}\\

\section*{Entradas}
\subsection*{Condiciones}
\begin{itemize}
\item Los alumnos tienen los conocimientos de programación suficientes para
  desarrollar un producto de software.
\end{itemize}
\subsection*{Productos de trabajo}
\begin{itemize}
\item \href{https://drive.google.com/open?id=18vx7IkGuZTE9kcgVFnWUrrSildkzyM8Y}{Planteamiento de Necesidades:}
\item \href{https://drive.google.com/open?id=1V-6SAoTrtV-phCGOPW19JJ2T-uBJpByW}{Guión del curso:}
\item Plantillas
\end{itemize}

\rule{0.8\textwidth}{.8pt}\\

\section*{Resultados}
\subsection*{Condiciones}
\begin{itemize}
\item Los alumnos han practicado conceptos básicos de Ingeniería de Software
  desarrollando un proyecto de software.
\end{itemize}

\subsection*{Productos de trabajo}
\begin{itemize}
\item Producto de software y su documentación.
\end{itemize}

\rule{0.8\textwidth}{.8pt}\\

\section*{Método} Método inicial de Desarrollo de Software.


\rule{0.8\textwidth}{.8pt}\\

\section*{Periodo}
Fecha de inicio: 5 de Agosto de 2019\\
\indent Fecha de fin: 5 de Diciembre de 2019\\

\rule{0.8\textwidth}{.8pt}\\

\section*{Reuniones del equipo}

\begin{center}
  \begin{table}[H]
  \centering
  \begin{tabular}{| c | c | c | c | c | c | c | c | }
    \hline
    hora & Lunes & Martes & Miércoles & Jueves & Viernes & Sábado & Domingo \\
    \hline
    14:00 & & & & & Presencial & & \\ \hline
    16:00-18:00 & & & & &  & Virtual & Virtual\\ \hline    
    19:00 & Presencial & Presencial & Presencial & Presencial &  &  &  \\ \hline
  \end{tabular}
  \caption{Horarios de reunión.}
  \label{tabla:horarios}
  \end{table}
\end{center}

\rule{0.8\textwidth}{.8pt}\\

\section*{Repositorio común de documentos}
\href{https://drive.google.com/drive/folders/18_KMcs7qbq7PF_Ij9WvfzrCTSS1qXZC6}{IS201-Ilakech-Alakech}
\end{document}
